% !TeX spellcheck = it_IT

\documentclass[handout]{beamer}

\usepackage[utf8]{inputenc}
\usepackage[italian]{babel}
\usepackage{multicol}
\usetheme{metropolis}

\usepackage{mathtools}
\DeclarePairedDelimiter{\ceil}{\lceil}{\rceil}


\title{Informatica B - Esercitazione 0}
\subtitle{Codifica binaria}

\author{Stefano Cereda\\
		stefano1.cereda@mail.polimi.it
	}
\date{26/09/2017}
\institute[PoliMi]{\vspace{0.5cm}\centering Politecnico di Milano \\ \vspace{0.2cm}
	\includegraphics[width=0.2\linewidth]{./logopolimi}}
%\logo{\includegraphics[width=15mm]{./logopolimi}}

%\setbeamercovered{invisible}

\makeindex

\begin{document}
\begin{frame}
	\maketitle
\end{frame}

\section{Introduzione}
\begin{frame}{Esercitazioni}
\begin{itemize}[<+->]
	\item stefano1.cereda@mail.polimi.it
	\item Esercitazioni = esercizi simili all'esame (teoria, esempi, esercizi)
	\item Approccio = teoria, esercitazioni, laboratorio
	\item Partecipare
	\item Consiglio carta e penna
	\item Dubbi: lezione, forum, mail, appuntamento
	\item Materiale online dopo esercitazioni
	\item Chi programma già?	
\end{itemize}
\end{frame}

\section{Alfabeti e codifiche}
\begin{frame}
\frametitle{Alfabeto}
\begin{itemize}[<+->]
	\item Alfabeto = insieme di simboli utilizzato (\alert{\{0,1\}} nella codifica binaria)
	\item L'informazione viene codificata da una sequenza di simboli (detti \emph{bit} nella codifica binaria)
	\item 8 \emph{bit} formano un \emph{byte}. I \emph{byte} vengono a loro volta raggruppati in \emph{word}.
\end{itemize}
\end{frame}

\begin{frame}
\frametitle{Alfabeti: numero di sequenze possibili}
Dato un alfabeto con 2 simboli, quante possibili sequenze di lunghezza \emph{N} possiamo rappresentare?
\pause

\iffalse
%\begin{table}[h]
%\begin{tabular}{cc}
%	\emph{N=2}&\\
%	0&0\\
%	0&1\\
%	1&0\\
%	1&1\\
%\end{tabular}
%\end{table}

%\begin{table}[h]
%\begin{tabular}{ccc}
%	\emph{N=3}&\\
%	0&0&0\\
%	0&0&1\\
%	0&1&0\\
%	0&1&1\\
%	1&0&0\\
%	1&0&1\\
%	1&1&0\\
%	1&1&1
%\end{tabular}
%\end{table}

\fi 

Le possibili combinazioni binarie di lunghezza \emph{N} sono $M = 2^N$

La lunghezza minima per rappresentare \emph{M} combinazioni è $N = \ceil{\log_2M}$ (arrotondamento per eccesso)
\end{frame}

\begin{frame}
\frametitle{Codifica}
\begin{itemize}[<+->]
	\item Il calcolatore deve rappresentare diversi tipi di dati.
	\item Codifica = \alert{regola} per mettere in corrispondenza i \emph{valori del dato} da codificare con una \emph{sequenza di simboli} dell'alfabeto.
	\item È una regola arbitraria che deve essere nota e rispettata da chi genere e da chi utilizza le informazioni (\emph{standard}).
	\item L'arbitrarietà consente di ottenere particolari \emph{proprietà}, decise in base all'uso che si farà delle informazioni.
\end{itemize}
\end{frame}

\section{Codifica binaria di numeri ed enumerazioni}

\subsection{Numeri naturali}
\begin{frame}
\frametitle{Codifica binaria di numeri naturali}
Questa codifica (\emph{binaria pura}) è simile a quella di tutti i giorni (posizionale in base 10).

\pause
\vskip .7cm
$N_b = a_n\cdot b^n + a_{n-1} \cdot b^{n-1} \dots a_1 \cdot p^1 + a_0 \cdot p^0 = \sum_{i=0}^n {a_i \cdot b^i}$

\pause\vskip .5cm
$1234_{10} = 1\cdot10^3 + 2\cdot10^2 + 3\cdot10^1 + 4\cdot10^0 = 1234_{10}$
\pause\vskip.3cm
$1010_{2} = 1\cdot2^3 + 0\cdot2^2 + 1\cdot2^1 + 0\cdot2^0 = 10_{10}$
\pause\vskip.3cm
$147_8 = 1\cdot8^2 + 4\cdot8^1 + 7\cdot8^0 = 103_{10}$
\pause\vskip.3cm
$B2A_{16} = 11\cdot16^2 + 2\cdot16^1 + 10\cdot16^0 = 2858_{10} $

\end{frame}

\begin{frame}
\frametitle{Conversione da decimale a binario}
Per convertire un numero decimale in binario puro, si utilizza il metodo delle \alert{divisioni successive}. Si divide ripetutamente il numero per 2 fino ad arrivare a zero. La codifica binaria è data dai resti delle divisioni (letti al contrario).

\begin{table}[h]
	\begin{tabular}{c|c}
		123 & 1 \\
		61  & 1 \\
		30  & 0 \\
		15  & 1 \\
		7   & 1 \\
		3   & 1 \\
		1	& 1	\\
		0   &  \\
	\end{tabular}
\end{table}
$123_{10} = 1111011_2$
\end{frame}

\begin{frame}
\frametitle{Esercizi conversioni decimale-binario}
Da fare:

$78_{10} = ?_{2}$

$135_{10} = ?_{2}$

$296_{10} = ?_{2}$

$17_{10} = 1000?_{2}$

$16_{10} = 1000?_{2}$

\end{frame}

\begin{frame}
\frametitle{Esercizi conversioni decimale-binario}
Risolti:

$78_{10} = 1001110_{2} = 1*2+1*4+1*8+1*64 = 78_{10}$

$135_{10} = 10000111_{2} = 1+2+4+128 = 135_{10}$

$296_{10} = 100101000_{2} = 8+32+256 = 296_{10}$

$17_{10} = 10001_{2}$

$16_{10} = 10000_{2}$

\end{frame}

\begin{frame}
\frametitle{Operazioni aritmetiche in binario puro}
\begin{table}[h]
	\begin{tabular}{c|c}
		0+0 & =0 \\
		0+1 & =1 \\
		1+0 & =1 \\
		1+1 & =0 con riporto di 1
	\end{tabular}
\end{table}
\begin{table}[h]
	\begin{tabular}{c|c}
		0-0 & =0 \\
		1-0 & =1 \\
		1-1 & =0 \\
		0-1 & =1 con prestito di 1
	\end{tabular}
\end{table}
\end{frame}

\begin{frame}{Somme, riporti, overflow e padding}
\begin{itemize}
	\item 10 + 01 = 11
	\item 01 + 01 = 10
	\item 11 + 01 = 00 con overflow
	\item 011 + 01 $\rightarrow$ 011 + 001 = 100
\end{itemize}
\end{frame}

\begin{frame}{Esercizi somma}
Da fare:

1010 + 0011 = ?

1011 + 0001 = ?

1011 + 1001 = ?

0100 + 111 = ?
\end{frame}

\begin{frame}{Esercizi somma}
Fatti:

1010 + 0011 = 1101 (10+3=13)

1011 + 0001 = 1100 (11+1=12)

1011 + 1001 = 0100 con overflow (11+10=8 ovf)

0100 + 111 = 1011 (4+7=11)
\end{frame}


\subsection{Enumerazioni}
\begin{frame}
\frametitle{Rappresentazione di enumerazioni}
Tramite una opportuna codifica possiamo rappresentare delle enumerazioni tramite cifre binarie.
\pause

Ad esempio, consideriamo l'insieme dei giorni della settimana. Dati $M=7$ valori distinti da rappresentare, abbiamo bisogno di $N=\ceil{\log_2 7}=\ceil{2.83}=3$ bit. 
\pause

\begin{table}[h]
	\begin{tabular}{cccc}
		lun&0&0&0\\
		mar&0&0&1\\
		mer&0&1&0\\
		gio&0&1&1\\
		ven&1&0&0\\
		sab&1&0&1\\
		dom&1&1&0
	\end{tabular}
\end{table}

Rimane una sequenza non utilizzata (111).
\end{frame}

\begin{frame}
\frametitle{ASCII}
Una delle codifiche per enumerazione più importanti è quella ASCII.

\href{https://it.wikipedia.org/wiki/ASCII}{https://it.wikipedia.org/wiki/ASCII} (notare la colonna CEC)

\begin{itemize}
	\item 128 caratteri (256 ASCII esteso)
	\item Caratteri di \emph{comando}, \emph{alfanumerici} e \emph{simboli}
	\item Proprietà interessanti
\end{itemize}

\end{frame}

\subsection{Numeri interi}
\begin{frame}
\frametitle{Numeri interi}
Settimana prossima!
\end{frame}

\end{document}
